% Shared experience block macros — used by resume.tex and postdoc-cv.tex.
% Each macro calls \cvitem{date}{content}, so the layout adapts to whichever
% \cvitem definition is active in the including document.

\newcommand{\expPostdocXray}{%
\cvitem{Jan 2025 -- Present}{%
\textbf{Postdoctoral Research Associate} \textit{(X-ray Scattering Signal Modeling)}\\%
\textit{Northeastern University, Dept. of Bioengineering,} Boston, MA\\[0.1em]%
\begin{highlights}%
\item Developed a correlation-aware feature reduction method to extract weak biological signatures from noisy X-ray scattering data, reducing 211 candidate features to 11 disease-discriminative markers.%
\item Derived Bayes-risk and approximation-error bounds that determine pruning limits for highly correlated data while preserving discriminative signal.%
\item Trained a compact, interpretable neural network with 174 parameters to estimate pathological protein aggregation, achieving 84.3\% test F1-score on held-out data.%
\end{highlights}%
}%
}

\newcommand{\expVesselSeg}{%
\cvitem{Apr 2026 -- Present}{%
\textbf{Computational Imaging Collaborator} \textit{(Neurovascular Segmentation)}\\%
\textit{Northeastern University, Dept. of Bioengineering,} Boston, MA\\[0.1em]%
\begin{highlights}%
\item Created an end-to-end training and inference pipeline to segment blood-vessel structures in noisy 3D fluorescence microscopy images, handling weak contrast, uneven focus, broken structures, and limited annotated data.%
\item Fine-tuned an Attention U-Net with a pretrained ResNet34 encoder, using attention-gated skip connections to infer latent vessel structure, guided by engineered image features and augmentation.%
\item Trained and validated the model on an HPC cluster with 5-fold cross-validation, checkpointing, W\&B tracking, attention-map review, and TP/FP/FN failure analysis, achieving approximately 91\% Dice on held-out test data.%
\end{highlights}%
}%
}

\newcommand{\expCOPD}{%
\cvitem{Nov 2025 -- Present}{%
\textbf{Ancillary Co-Principal Investigator} \textit{(Multi-Omics Biomarkers of Lung Disease)}\\%
\textit{Brigham and Women's Hospital / Harvard Medical School,} Boston, MA\\[0.1em]%
\begin{highlights}%
\item Lead transcriptomic analysis of COPDGene participants ($n \approx 3{,}000$) to identify molecular signatures associated with mechanically affected lung (MAL) using participant stratification, differential gene expression, and pathway enrichment analysis.%
\item Develop proteomics mediation models to identify plasma proteins that mediate the relationship between mechanically affected lung (MAL) and 5-year emphysema progression.%
\end{highlights}%
}%
}

\newcommand{\expPhDXray}{%
\cvitem{Oct 2020 -- Dec 2024}{%
\textbf{Graduate Research Assistant} \textit{(Molecular Mapping of Alzheimer's Pathology)}\\%
\textit{Northeastern University, Dept. of Electrical \& Computer Engineering,} Boston, MA\\[0.1em]%
\begin{highlights}%
\item Designed an inverse-problem framework that generates Alzheimer's pathology maps from raw X-ray microdiffraction data, partnering with wet-lab teams at MGH for histology and biophysicists at Brookhaven National Laboratory for synchrotron measurements across 15 human brain samples.
\item Formulated a pipeline that chains background subtraction, t-SNE, Gaussian Mixture Models, and a neural network to generate probabilistic spatial maps of pathological protein aggregates from noisy X-ray scattering measurements.%
\item Identified a consistent molecular signature of Alzheimer's cross-$\beta$ fibrils: a $\sim$4.7~\AA{} $\beta$-strand spacing with an approximately $-1^\circ$ left-handed twist that persisted across disease stages.%
\item Authored 5 peer-reviewed publications and delivered 5 conference presentations, including 2 invited talks recognized with student lecturer awards.%
\end{highlights}%
}%
}

\newcommand{\expCryoET}{%
\cvitem{Sep 2022 -- Dec 2022}{%
\textbf{Visiting Graduate Researcher} \textit{(Cryo-ET Automation)}\\%
\textit{Harvard Medical School, Dept. of Biological Chemistry \& Molecular Pharmacology,} Boston, MA\\[0.1em]%
\begin{highlights}%
\item Designed a Python-based GUI to automate cryo-electron tomography data acquisition, including microscope setup, motion correction, drift correction, and tilt-series alignment.%
\item Reduced setup time by 40\% and increased throughput for structural biology workflows at Boston-area hospitals including Boston Children's Hospital, Dana-Farber, and MGH.%
\end{highlights}%
}%
}

\newcommand{\expDrones}{%
\cvitem{Jul 2019 -- Aug 2020}{%
\textbf{Graduate Research Assistant} \textit{(Drone Swarm Sensing and Control)}\\%
\textit{Northeastern University, Dept. of Electrical \& Computer Engineering,} Boston, MA\\[0.1em]%
\begin{highlights}%
\item Engineered custom autonomous drones that combined camera, infrared, radar, and GPS data to detect and locate people and vehicles in outdoor environments with limited wireless communication.%
\item Implemented drone swarm behaviors for coordinated flight, target tracking, and collision avoidance, demonstrated to U.S. Air Force leadership at Unmanned Aerial Systems Pitch Day.%
\end{highlights}%
}%
}

\newcommand{\expToad}{%
\cvitem{Jan 2018 -- Jun 2019}{%
\textbf{Graduate Research Assistant} \textit{(Wildlife Signal Detection and Sensing Device)}\\%
\textit{Texas State University, Dept. of Electrical Engineering,} San Marcos, TX\\[0.1em]%
\begin{highlights}%
\item Prototyped and deployed a solar-powered Raspberry Pi acoustic monitor with onboard inference and cellular alerting for endangered Houston toad detection, scaling the system to 50 field units for conservation use in Bastrop County, Texas.%
\item Implemented power spectral density analysis, tuned bandpass filtering, MFCC extraction, and an LSTM classifier to distinguish Houston toad mating calls from other animal sounds and field noise, achieving 98\% true-positive and 92.6\% test accuracy on field audio.%
\end{highlights}%
}%
}
